\chapter{Dialogo Testuale con le
Tradizioni}\label{dialogo-testuale-con-le-tradizioni}

\begin{quote}
\textbf{Scopo dell'appendice.} Espandere la mappa concettuale del Cap.
13.3 in dieci dialoghi brevi e diretti con le tradizioni che la DEQ$^2$
eredita, modifica o supera. Ciascun dialogo è strutturato come
\emph{obiezione attesa} dalla tradizione + \emph{risposta DEQ$^2$} +
\emph{concessione esplicita} di ciò che la DEQ$^2$ \emph{non} fa rispetto a
quella tradizione. La forma è quella del dialogo immaginario: non
pretende di parlare per gli autori citati, pretende di \emph{farsi
parlare} dalle loro categorie.
\end{quote}



\section{\texorpdfstring{Con Giovanni Arrighi (\emph{Il lungo XX
secolo},
1994)}{Con Giovanni Arrighi (Il lungo XX secolo, 1994)}}\label{con-giovanni-arrighi-il-lungo-xx-secolo-1994}

\textbf{Arrighi:} ``Lei ha individuato una transizione 2028-2031 senza
riconoscere che è la \emph{fase autunnale} del ciclo finanziario USA ---
già descritta nel mio ciclo egemonico
genovese-olandese-britannico-americano. Non sta scoprendo nulla, sta
misurando ciò che io ho narrato.''

\textbf{DEQ$^2$:} ``La concessione è piena: Lei ha narrato la \emph{forma}
della transizione; io provo a \emph{misurarne il bordo}. La
\(T_s^{(2),\text{sys}} \to 0\) del Cap. 7 sull'USA è la traduzione
formale del Suo \emph{signal crisis} (anni '70-'80) seguito da
\emph{terminal crisis} (anni 2020-2030). Quello che aggiungo: la
transizione \emph{non genera automaticamente} un nuovo egemone --- può
generare invece un'asimmetria multi-attore in cui \(\Phi\) esportabile
risiede in un Bilanciere (Brasile) anziché in un nuovo egemone (Cina).
Lei pensava in termini di \emph{successione}; io penso in termini di
\emph{biforcazione}.''

\textbf{Concessione.} Senza Arrighi, il Cap. 7 e il Cap. 13.7 sarebbero
illeggibili. La DEQ$^2$ \emph{eredita} l'orizzonte temporale e l'ipotesi di
transizione; \emph{aggiunge} la quantificazione e la possibilità di un
esito non-egemonico.



\section{\texorpdfstring{Con Peter Turchin (\emph{Secular Cycles}, 2009;
\emph{End Times},
2023)}{Con Peter Turchin (Secular Cycles, 2009; End Times, 2023)}}\label{con-peter-turchin-secular-cycles-2009-end-times-2023}

\textbf{Turchin:} ``La Sua \(\Omega\) è la mia \emph{elite
overproduction}; il Suo \(\varepsilon_{\text{dis}}\) è la mia
\emph{Popular Immiseration}. Ha riformulato in notazione greca quello
che ho già modellato cliodinamicamente.''

\textbf{DEQ$^2$:} ``Ha ragione su due variabili su quattro. Ma manca
\(\Phi\). Lei modella il \emph{denominatore} della stabilità sistemica
(involuzione + immiserazione \(\Rightarrow\) instabilità) ma non il
\emph{numeratore} (coerenza di fase \(\times\) antifragilità
\(\Rightarrow\) stabilità positiva). Senza \(\Phi\), la cliodinamica
predice \emph{quando} un sistema crolla, ma non \emph{cosa lo riempie}.
La differenza è cruciale per il 2031: due sistemi possono avere la
stessa \(\Omega + \varepsilon_{\text{dis}}\) e finire in M1 (decoerenza
con \(\Phi\)-engine) o in M2 (decoerenza degenere). La cliodinamica è
agnostica fra M1 e M2; la DEQ$^2$ no.''

\textbf{Concessione.} Le serie temporali storiche di Turchin sono
essenziali per calibrare \(A\) via stress-test retrospettivo (App. C.3).
La DEQ$^2$ rinuncia a fare previsioni a tre secoli (orizzonte cliodinamico)
--- si limita a un \emph{ciclo} (2026-2031).



\section{\texorpdfstring{Con Immanuel Wallerstein (\emph{The Modern
World-System},
1974-2011)}{Con Immanuel Wallerstein (The Modern World-System, 1974-2011)}}\label{con-immanuel-wallerstein-the-modern-world-system-1974-2011}

\textbf{Wallerstein:} ``Il Brasile è semi-periferia ascendente del
sistema-mondo capitalista. Lei lo presenta come \(\Phi\)-engine --- cioè
come un attore con prerogative \emph{centrali} di produzione di ordine.
È un'illusione: la semi-periferia \emph{non} genera coerenza globale, è
generata da essa.''

\textbf{DEQ$^2$:} ``Concedo che lo schema centro-periferia ha potere
descrittivo notevole per il periodo 1500-2000. Ma la sua linearità
(centro \(\to\) semi-periferia \(\to\) periferia) è un \emph{vincolo}
che la DEQ$^2$ rifiuta di importare. La \(\Phi\) non è correlata alla
posizione nel sistema-mondo: dipende dalla struttura interna del sistema
sociale di un paese (Cap. 4.5, Cap. 5.2). Il Brasile può essere
\emph{semi-periferico per output} e \emph{centrale per \(\Phi\)
esportabile} --- è esattamente la posizione strutturale del Cap. 9. Lei
vincola; io disvincolo.''

\textbf{Concessione.} Lo schema centro-periferia resta utile per
descrivere il \emph{flusso materiale} (capitale, manodopera, risorse).
Ma il \emph{flusso di coerenza} segue regole diverse --- l'errore
strutturale del wallersteinismo è confondere i due flussi.



\section{\texorpdfstring{Con Amado Cervo (\emph{Inserção internacional},
2008)}{Con Amado Cervo (Inserção internacional, 2008)}}\label{con-amado-cervo-inseruxe7uxe3o-internacional-2008}

\textbf{Cervo:} ``L'Estado Logístico è una categoria che ho costruito
per descrivere il Brasile post-2003. Lei lo importa, ma la matematica
DEQ$^2$ rischia di trasformare una \emph{categoria politica vivente} in un
\emph{parametro statico}. La politica estera brasiliana non è un
coefficiente.''

\textbf{DEQ$^2$:} ``Concessione totale. La
\(\Pi^{\text{Logístico}} = \sum_k M_e^{(k)} Q_0^{(k)} (1 - \delta_R^{(k)} z_R^{(k)})\)
del Cap. 9.1 \emph{non sostituisce} la Sua categoria --- la
\emph{traduce} in linguaggio operativo per la dashboard SPC-DEQ. La
traduzione è utile \emph{operativamente} (permette monitoraggio in tempo
reale) ma è \emph{parziale concettualmente} (perde la dimensione storica
accumulativa che Lei evidenzia). Per questo l'Estado Logístico nella
DEQ$^2$ è \emph{condizionale} alla continuità politica --- esattamente la
fragilità del Cap. 9.5 + Cap. 12.5.''

\textbf{Concessione.} La DEQ$^2$ rinuncia a teorizzare il \emph{processo
storico} di formazione dell'Estado Logístico (1990-2003). Si limita a
operativizzarne i parametri \emph{date le condizioni del 2025-2031}.



\section{\texorpdfstring{Con Nassim Nicholas Taleb (\emph{Antifragile},
2012)}{Con Nassim Nicholas Taleb (Antifragile, 2012)}}\label{con-nassim-nicholas-taleb-antifragile-2012}

\textbf{Taleb:} ``Lei prende la mia antifragilità individuale e la
incolla a un sistema egemonico globale. La mia teoria è
\emph{anti-sistemica}: l'antifragilità \emph{vive nei dettagli e nelle
code}, non nelle medie aggregate. Il Suo \(\langle A\rangle\) è
esattamente l'errore che combatto.''

\textbf{DEQ$^2$:} ``Concessione parziale. Sì, \(\langle A\rangle\) è una
media --- ma è modulata da \(\Phi\), e l'asimmetria
\(\Phi\)-moltiplica-\(A\) vs \(\Phi\)-divide-\(\Omega\) è esattamente la
firma matematica della Sua intuizione: il sistema \emph{amplifica} la
fragilità più di quanto amplifichi l'antifragilità. La DEQ$^2$ formalizza
la convessità che Lei ha portato come euristica. Inoltre il Cap. 12.6
(\(S_{\mathcal{R}_-}\)) è interamente \emph{taleb-iano}: opzionalità,
ridondanza, disaccoppiamento dalla traiettoria sistemica. La DEQ$^2$
accoglie la Sua critica trasformandola in protocollo.''

\textbf{Concessione.} La DEQ$^2$ \emph{non} è una teoria della
\emph{previsione} nel senso strong (rifiutato da Taleb); è una teoria
dello \emph{spazio degli scenari} con probabilità soggettive esplicite.
Sotto questo profilo è coerente con \emph{Black Swan} e \emph{Skin in
the Game}.



\section{\texorpdfstring{Con Antonio Gramsci (\emph{Quaderni del
carcere}, 1929-1935) e Robert Cox (\emph{Production, Power, and World
Order},
1987)}{Con Antonio Gramsci (Quaderni del carcere, 1929-1935) e Robert Cox (Production, Power, and World Order, 1987)}}\label{con-antonio-gramsci-quaderni-del-carcere-1929-1935-e-robert-cox-production-power-and-world-order-1987}

\textbf{Gramsci/Cox:} ``L'egemonia non è \(T_s^{(2),\text{sys}}\). È
rapporto fra forze materiali, istituzionali e culturali in cui il
consenso e la coercizione si combinano \emph{organicamente}. Lei riduce
l'organico al numerico.''

\textbf{DEQ$^2$:} ``Concedo che \(T_s^{(2),\text{sys}}\) è una
\emph{misura} dell'egemonia, non la \emph{spiegazione} dell'egemonia. La
spiegazione resta gramsciana: il consenso (\(s_{iii}\) DEQ$^2$) e la
coercizione (\(\sigma \cdot \varepsilon_{\text{dis}}\)) si combinano in
proporzioni che variano per blocco storico. Quello che la DEQ$^2$ aggiunge
è la \emph{condizione di stabilità}: per qualunque blocco storico,
\(T_s^{(2),\text{sys}} > 1\) è condizione necessaria di durabilità. La
DEQ$^2$ è dunque \emph{interna} alla teoria gramsciana dell'egemonia, non
alternativa.''

\textbf{Concessione.} La nozione di \emph{blocco storico} gramsciano
contiene una temporalità \emph{qualitativa} (cesura, conservazione,
rivoluzione passiva, restaurazione) che la DEQ$^2$ non cattura. La
transizione di fase del Cap. 3.5.3 corrisponde solo alla \emph{cesura}
gramsciana; le altre forme richiedono estensione.



\section{\texorpdfstring{Con Edward Luttwak (\emph{Strategy: The Logic
of War and Peace},
1987-2001)}{Con Edward Luttwak (Strategy: The Logic of War and Peace, 1987-2001)}}\label{con-edward-luttwak-strategy-the-logic-of-war-and-peace-1987-2001}

\textbf{Luttwak:} ``La grande strategia è \emph{paradossale}: ciò che
funziona si esaurisce, ciò che non funziona viene ripreso. La Sua
matematica è \emph{lineare} nel tempo. Si perde il paradosso.''

\textbf{DEQ$^2$:} ``Concessione. Le forme \(T_s^{(2),\text{sys}}\) sono
monotoniche (Cap. 3.5.1); la dinamica del Cap. 5.4 ammette però
retroazioni non lineari (R6 \(\Omega \to \varepsilon_{\text{dis}}\), R3
\(\Phi \to \varepsilon_{\text{dis}}\) negativa) che producono
comportamento \emph{paradossale} nell'evoluzione di
\(\boldsymbol{p}(t)\): la stessa azione può essere antifragile in M1 e
fragile in M2. È il paradosso luttwakiano nella forma DEQ$^2$: la strategia
ottima dipende dallo scenario realizzato, e lo scenario realizzato non è
noto al momento della scelta.''

\textbf{Concessione.} La DEQ$^2$ \emph{non} analizza il \emph{registro
militare} della grande strategia (manovra, deception, logistica). Si
limita al \emph{registro strutturale} (coerenza, fragilità,
involuzione). Per la grande strategia operativa Lei resta riferimento.''



\section{Con la Scuola di Francoforte (Marcuse, Adorno, Habermas) ---
via
Topografia}\label{con-la-scuola-di-francoforte-marcuse-adorno-habermas-via-topografia}

\textbf{Marcuse/Adorno/Habermas:} ``L'apparato critico francofortese ha
denunciato la \emph{razionalità strumentale} come distruttrice della
razionalità sostanziale. Lei produce un'altra razionalità strumentale:
la matematizzazione del campo discorsivo. La denuncia, applicata a Lei
stesso, è devastante.''

\textbf{DEQ$^2$:} ``Concedo l'auto-implicazione. Difesa: la formalizzazione
del campo \(\Pi\) topografico (eredità Topografia \S{}5) \emph{non} riduce
\(z\) --- al contrario, \emph{misura quanto \(z\) è presente}. La
\(T_s^{(2),\text{sys}}\) premia sistemi con \(z\) alto (asse della
complessità dialettica al numeratore di \(Q_0\)). Una razionalità
strumentale che misura --- e premia --- la razionalità sostanziale è
\emph{forse} l'unica via per parlare alla decisione politica nel 2026
senza essere tagliati fuori dal discorso pubblico. Habermas avrebbe
odiato la DEQ$^2$; ma Habermas pubblicava libri come
\emph{Legitimationsprobleme} destinati ad accademici, non a policymaker.
Il Bilanciere parla alla decisione.''

\textbf{Concessione.} La DEQ$^2$ \emph{non} è critica francofortese --- è
\emph{sintesi tecnica della critica francofortese}. Perde
l'incandescenza dialettica originaria; guadagna in azionabilità.
Trade-off accettato.



\section{\texorpdfstring{Con Xiang Biao (\emph{Self as Method}, 2020;
\emph{Anti-involution},
2021-2024)}{Con Xiang Biao (Self as Method, 2020; Anti-involution, 2021-2024)}}\label{con-xiang-biao-self-as-method-2020-anti-involution-2021-2024}

\textbf{Xiang Biao:} ``Il \emph{Neijuan} è una categoria etnografica
viva, nata nei dormitori delle università cinesi negli anni 2010 e
diffusa nei contesti 996. Lei la trasforma in \(\Omega\) --- un
parametro di equazione. Le persone che vivono il Neijuan non vivono
\(\Omega\); vivono fatica.''

\textbf{DEQ$^2$:} ``Concessione totale. Il passaggio dal Neijuan
etnografato a \(\Omega\) misurato è una \emph{perdita di contenuto
fenomenologico}. Difesa: la perdita è funzionale a un obiettivo diverso
dal Suo. La Sua etnografia \emph{fa sentire} il problema; la DEQ$^2$
\emph{rende contabile} il problema. Le due operazioni sono
complementari, non sostitutive. Inoltre il Cap. 8.2 cita Lei
\emph{direttamente} come \emph{firma empirica} della variabile --- non
come decorazione retorica. Senza \emph{Self as Method}, il Cap. 8
sarebbe un esercizio aggregato senza ancoraggio empirico.''

\textbf{Concessione.} La DEQ$^2$ \emph{non} può sostituire la ricerca
etnografica. Le serve come \emph{lettore di senso}; senza, leggerebbe
deflazione PPI senza vedere persone.''



\section{\texorpdfstring{Con Yoshiki Kuramoto (\emph{Chemical
Oscillations, Waves, and Turbulence},
1984)}{Con Yoshiki Kuramoto (Chemical Oscillations, Waves, and Turbulence, 1984)}}\label{con-yoshiki-kuramoto-chemical-oscillations-waves-and-turbulence-1984}

\textbf{Kuramoto:} ``Il mio modello descrive oscillatori chimico-fisici.
Voi sociologi lo applicate da decenni a sistemi sociali con disinvoltura
preoccupante. Le ipotesi (campo medio, simmetria della distribuzione
\(g(\omega)\), accoppiamento sinusoidale) \emph{non} sono verificabili
sui sistemi sociali.''

\textbf{DEQ$^2$:} ``Concessione: l'Appendice A.2 dichiara esplicitamente A1
(Kuramoto applicabile) come \emph{ipotesi operativa}, non come fatto
verificato. Difesa: la matematica di Kuramoto è la \emph{più semplice}
descrizione di sincronizzazione fra oscillatori non identici. Per \(N\)
grande e accoppiamenti deboli, qualunque modello di sincronizzazione
converge in qualche limite a Kuramoto (universalità). La DEQ$^2$ è dunque
\emph{minimalista}: usa la più povera matematica di sincronizzazione
disponibile, e ne dichiara i limiti. Estensioni a Kuramoto multi-cluster
(Strogatz 2003) sono pianificate per il paper EN --- proprio per i
fenomeni che il modello base non cattura, come la polarizzazione USA.''

\textbf{Concessione.} La DEQ$^2$ \emph{non} dimostra che Kuramoto è la
matematica corretta del sociale. Argomenta che è una \emph{prima
approssimazione consistente}. La verifica empirica è demandata al paper
EN --- coerente con la posizione popperiana del Cap. 11.''



\section{Sintesi dei Dieci Dialoghi}\label{sintesi-dei-dieci-dialoghi}

\Proved{} \textbf{Pattern emergente.} In ciascuno dei dieci dialoghi:

\begin{itemize}
\tightlist
\item
  la DEQ$^2$ fa una \textbf{concessione esplicita} (cosa \emph{non} fa
  rispetto alla tradizione);
\item
  la DEQ$^2$ fa una \textbf{risposta sostanziale} (cosa \emph{aggiunge}
  alla tradizione);
\item
  la DEQ$^2$ fa una \textbf{clausola di apertura} (estensioni / verifiche
  pianificate per il futuro).
\end{itemize}

\Hypoth{} \textbf{Posizione metodologica.} La DEQ$^2$ non è teoria
\emph{contro}; è teoria \emph{con}. Questo è coerente con l'ipotesi di
campo medio: ogni teoria è un \emph{oscillatore} nel campo
intellettuale; la coerenza di fase del campo
\(\Phi^{\text{intellettuale}}\) è promossa da chi accoglie e modifica,
non da chi denuncia e sostituisce. La forma stessa di questa appendice è
\emph{applicazione riflessiva} della DEQ$^2$ alla propria posizione nel
campo del sapere.

\Proved{} \textbf{Conseguenza editoriale.} I lettori specialisti di una
qualunque delle dieci tradizioni elencate trovano nella DEQ$^2$: 1.
\emph{Riconoscimento esplicito} della loro tradizione (cosa accolto); 2.
\emph{Argomentazione esplicita} del valore aggiunto (cosa nuovo); 3.
\emph{Limite esplicito} del valore aggiunto (cosa lasciato a loro).



\section{Tradizioni Non Dialogate}\label{tradizioni-non-dialogate}

\Conj{} \textbf{Esclusioni esplicite.} Per ragioni di spazio, la DEQ$^2$
\emph{non} dialoga in questa appendice con:

\begin{itemize}
\tightlist
\item
  \textbf{Foucault} (potere capillare, biopolitica) --- il dialogo
  richiederebbe una riformulazione del campo \(\Pi\) in termini di
  \emph{dispositivo} anziché di \emph{accoppiamento}. Rinviato a futura
  estensione.
\item
  \textbf{Latour} (Actor-Network Theory) --- la DEQ$^2$ assume attori
  discreti; ANT li rifiuta. Trade-off metodologico irrisolto.
\item
  \textbf{Bourdieu} (campo, habitus, capitale) --- il \emph{campo}
  bourdieusiano è strutturalmente affine al campo \(\Pi\) topografico,
  ma con \emph{abito disposizionale} anziché \emph{fase di Kuramoto}.
  Riformulazione possibile, non tentata qui.
\item
  \textbf{Tilly} (contentious politics) --- la dinamica dei movimenti
  sociali è esterna al perimetro DEQ$^2$; integrazione possibile via
  estensione \(\dot\Phi\) con termine di forzatura mobilitazionale.
\end{itemize}

\Proved{} \textbf{Posizione.} Le quattro esclusioni sopra non sono
dimissioni; sono \emph{aperture pianificate} per estensioni successive
del programma DEQ$^2$.



\section{§D.4 — Le Radici Evolutive: Hamilton, Wilson-Sober, Nowak,
Kropotkin}\label{d4-radici-evolutive}

\begin{quote}
\textbf{Scopo di questa sezione.} I quattro dialoghi che seguono
portano alla DEQ$^2$ la sua \emph{microfondazione causale evolutiva}: la
spiegazione di \emph{perché} $\Phi$ emerge, di \emph{come} si erode e di
\emph{cosa significa} la decoerenza alla scala degli organismi e delle
popolazioni. Non è un'aggiunta decorativa --- è il fondamento causale
della struttura formale. La completa derivazione del framework CPGG si
trova in \emph{Captured Cooperation and Phase Decoherence} (Marcelino
2026, working paper).
\end{quote}

\subsection{Con William D. Hamilton (\emph{The Genetical Evolution of
Social Behaviour}, 1964)}\label{con-hamilton-1964}

\textbf{Hamilton:} ``Lei usa il parametro d'ordine di Kuramoto per
descrivere la cooperazione sistemica. Ma da dove viene la cooperazione?
Lei la presuppone come fatto; io la derivo dalla genetica. Senza la
condizione $rB > C$, $\Phi > 0$ è una quantità orfana di causa.''

\textbf{DEQ$^2$:} ``È la critica più profonda che si potesse fare al
framework originale --- ed è corretta. La risposta sta nella
\emph{corrispondenza strutturale} (non identità) tra le sue condizioni
genetiche e le condizioni sistemiche. La condizione $rB > C$ si traduce
nel sistema geopolitico come: la cooperazione ($\Phi > 0$) è stabile
quando il beneficio collettivo della sincronia supera il costo
individuale di cooperare --- e questo è diventato, nella DEQ$^2$
aggiornata (Cap. 3, §3.2), la giustificazione causale di $\Phi$, non solo
la sua descrizione. Lei, Hamilton, ha fornito il \emph{perché} che io
non avevo.''

\textbf{Concessione.} La condizione $rB > C$ è qui applicata come
\emph{isomorfismo strutturale} (\Hypoth{}) a popolazioni di attori
geopolitici, non a popolazioni di organismi biologici. Il trasferimento
richiede l'ipotesi H1 (§4.0) che il campo $\Pi_j$ si mappa su una fase
--- un'ipotesi dichiarata, non un risultato derivato. La deriva di
Hamilton è preservata con onestà epistemica.

\subsection{Con David Sloan Wilson e Elliott Sober (\emph{Unto Others},
1994)}\label{con-wilson-sober-1994}

\textbf{Wilson \& Sober:} ``La letteratura ha a lungo trattato la
selezione di gruppo come irrilevante, seguendo Dawkins. Lei ha inizialmente
citato Dawkins a supporto del suo cambio di livello descrittivo --- un
errore che abbiamo visto molte volte. La \emph{selezione di gruppo} non
è un'estensione opzionale della teoria evolutiva: è la struttura che
giustifica \emph{qualsiasi} affermazione su benefici collettivi.''

\textbf{DEQ$^2$:} ``Avete pienamente ragione, e la prima bozza del
framework era difettosa esattamente su questo punto: Dawkins (1976)
rigetta la selezione di gruppo --- citarlo per giustificare il cambio di
livello micro/macro era contraddittorio. La correzione è implementata
(Cap. 3, §3.2): il fondamento formale è la vostra teoria della selezione
multilivello, formalizzata da Okasha (2006), combinata con l'equazione di
Price (1970) come identità algebrica esatta. L'equazione di Price
decompone $\Delta\bar{z}$ in termine inter-gruppo (proporzionale a $\Phi$)
e intra-gruppo (proporzionale a $-\Omega$) --- una struttura che giustifica
formalmente la forma del numeratore e del denominatore di
$T_s^{(2),\text{sys}}$.''

\textbf{Concessione.} La selezione multilivello DEQ$^2$ opera su sistemi
socio-istituzionali, non su alleli. L'applicazione richiede che il
``genotipo'' sia interpretato come configurazione istituzionale
(\Hypoth{}), e la ``fitness'' come payoff sistemico. Questa
reinterpretazione è dichiarata come ipotesi operativa, non come fatto
biologico.

\subsection{Con Martin A. Nowak (\emph{Five Rules for the Evolution of
Cooperation}, 2006)}\label{con-nowak-2006}

\textbf{Nowak:} ``Ha un problema di identificazione. I suoi attori ---
USA, Cina, Brasile, Musk --- presentano caratteristiche di tutti e cinque
i meccanismi cooperativi simultaneamente: reciprocità diretta (via trattati
bilaterali), reciprocità indiretta (via reputazione), selezione di rete
(via posizionamento geopolitico), selezione di gruppo (via alleanze),
selezione di parentela (via affinità culturale). Come distingue quale
meccanismo guida la dinamica?''

\textbf{DEQ$^2$:} ``La sua tassonomia dei cinque meccanismi diventa
nella DEQ$^2$ una \emph{tassonomia diagnostica degli attori}, non una
classificazione esclusiva dei meccanismi. Il CPGG (§D.4, working paper)
focalizza su un meccanismo specifico --- la \emph{cattura istituzionale}
che degrada tutti e cinque i meccanismi simultaneamente. La tesi è:
quando $\phi$ (la frazione estrattiva) supera $\phi^*$, tutti i cinque
meccanismi di Nowak vengono simultaneamente indeboliti --- la condizione
$b/c > k_{\text{eff}}(\phi)$ cessa di valere per ciascuno di essi. La
DEQ$^2$ non sceglie quale meccanismo; diagnostica quando tutti
collassano.''

\textbf{Concessione.} Il CPGG modella la cattura istituzionale, non la
dinamica evolutiva di ciascuno dei cinque meccanismi separatamente. Una
modellazione completa richiederebbe cinque CPGG separati o un modello
unificato di interazione tra meccanismi. Questo è F1 e F2 nell'agenda
formale aperta.

\subsection{Con Pyotr Kropotkin (\emph{Mutual Aid: A Factor of
Evolution}, 1902)}\label{con-kropotkin-1902}

\textbf{Kropotkin:} ``Ho documentato, specie dopo specie e cultura dopo
cultura, che la cooperazione è la norma e la competizione è l'eccezione.
Lei costruisce un formalismo per misurare la decoerenza --- che è
esattamente ciò che ho chiamato \emph{rottura del mutuo appoggio}. Ma
venti anni dopo, il Neijuan e la cattura estrattiva sembrano
schiaccianti. Ha dimenticato la mia tesi fondamentale?''

\textbf{DEQ$^2$:} ``No, l'ho \emph{rifondato}. La sua intuizione ---
$\Phi > 0$ come stato naturale, la decoerenza come stato prodotto --- è
ora formalizzata nella Proposizione K1 (§8 del framework CPGG): in
assenza di cattura istituzionale ($\sigma_{\text{san}} = \sigma_0$,
$\kappa_E = 0$), la frazione estrattiva di equilibrio $\phi_{\text{nat}}$
soddisfa $\phi_{\text{nat}} < \phi^*(\sigma_0)$ per parametri
biologicamente plausibili. La cooperazione è l'esito stabile; la
decoerenza richiede la condizione prodotta $\alpha\phi + \beta\kappa_E >
\log(\sigma_0/\sigma_{\text{crit}})$. Lei, Kropotkin, aveva torto
sull'inevitabilità della cooperazione nella storia --- non perché la
cooperazione sia rara, ma perché la cattura è più facile dell'apertura.
Il suo \emph{mutuo appoggio} è la condizione di default, non il
risultato garantito.''

\textbf{Concessione.} L'inversione epistemica --- decoerenza come stato
prodotto, cooperazione come stato naturale --- ha il rischio di diventare
normativamente ingenua: ``basata sulla natura, quindi dovuta''. La
\emph{naturalistica fallacy} è gestita fondando la parte normativa del
Cap. 12 sui principi di design di Ostrom (1990) piuttosto che su
Kropotkin direttamente. Ostrom fornisce un fondamento empirico: i commons
\emph{possono} essere gestiti cooperativamente, sotto condizioni
istituzionali precise --- non \emph{naturalmente} o \emph{inevitabilmente}.
Kropotkin è ispirazione fenomenologica; Ostrom è fondamento normativo.



\section{Sintesi dell'Appendice}\label{sintesi-dellappendice}

\Proved{} \textbf{Risultato.} La DEQ$^2$ si posiziona esplicitamente in
quattordici tradizioni intellettuali contemporanee (dieci originali +
quattro evolutive in §D.4), dichiarando per ciascuna concessione, risposta
e clausola di apertura. La forma del dialogo immaginario è
\emph{applicazione riflessiva} dell'apparato (campo medio Kuramoto
applicato al campo intellettuale). Quattro tradizioni ulteriori (Foucault,
Latour, Bourdieu, Tilly) sono dichiarate come estensioni pianificate.

\Hypoth{} \textbf{Tesi conclusiva del libro.} La DEQ$^2$ esiste \emph{fra}
tradizioni, non \emph{contro} di esse --- coerentemente con la sua tesi
sostanziale che la \(\Phi\) vera si genera per accoppiamento spontaneo
di sotto-attori, non per imposizione gerarchica. Il libro stesso è un
\emph{atto \(\Phi\)-engine} nel campo intellettuale: produce coerenza
ammettendo dissonanza, e si offre alla critica come condizione di vita.



\emph{Appendice di posizionamento. Stile: dialogo immaginario,
concessioni esplicite, aperture dichiarate.}
